\section{PuzzleNet: a Neural Model of Tactic and Difficulty}\label{sec:puzzlenet}

Two measured weaknesses motivate a learned model, and both are documented in
Chapter~\ref{ch:results}. First, every weakness estimate in the system inherits the tactic labeller's
errors, and the hand-written tagger of Section~\ref{sec:features} agrees with Lichess themes on only
58\% of held-out puzzles ($\kappa = 0.52$), while never emitting five of the categories at all. The
label-quality experiment (E3) shows that perfect labels would raise early targeting well beyond what
the tagger achieves, so labelling is a real bottleneck rather than a cosmetic one. Second, a mined
puzzle's difficulty is a closed-form guess from the size of the player's error and the length of the
solution; it is uncalibrated, carries no uncertainty, and 67.5\% of all real attempts were on mined
puzzles.

Both are supervised learning problems for which a large, free, labelled dataset already exists: the
5.88 million Lichess puzzles, each carrying theme tags and a Glicko-2 rating with its own deviation.
\textbf{PuzzleNet} is a single multi-task network trained on that dataset. It reads a position and the
line that follows it, and predicts the line's tactical category, its Lichess theme tags, and how
difficult it is, with an uncertainty attached to that difficulty.

\subsection{Task definition}

Let $s_0$ be a position with the \emph{solver} to move and $L = (m_1, r_1, m_2, r_2, \dots)$ the
principal line from it, alternating solver moves $m_i$ and replies $r_i$. The same pair $(s_0, L)$
occurs in all three places the system needs a label: a Lichess puzzle (the solution line), a critical
position in an analysed game (Stockfish's principal variation), and a mined puzzle (the forced
continuation built by the miner). The network estimates three things:

\begin{align}
  P(\text{category} = c \mid s_0, L), \quad c \in \{1,\dots,24\}, \\
  P(\text{theme}_j = 1 \mid s_0, L), \quad j \in \{1,\dots,66\}, \\
  b \mid s_0, L \;\sim\; \mathcal{N}\!\left(\mu(s_0, L),\; \sigma^2(s_0, L)\right),
\end{align}

\noindent where $b$ is the puzzle's difficulty on the Lichess rating scale. The theme head is an
auxiliary task: theme tags are a richer description of the same line, and predicting them gives the
shared trunk more to learn from than a single label per puzzle.

\subsection{Input representation}

The input is a fixed vector of 3{,}329 features. Its design follows two principles: give the network
the raw board, and also give it the relations a tactician would look for, so that it need not
rediscover the rules of attack and defence from piece positions alone. Whether the second part earns
its place is an ablation, not an assumption (Section~\ref{sec:ablations}).

Everything is encoded \emph{from the solver's point of view}. When Black is to move, squares are
mirrored ($a1 \leftrightarrow a8$) and colours swapped, so that the solver's pieces always occupy the
first six planes and the solver's pawns always advance up the board. A position and its
colour-mirrored twin therefore produce the identical vector, which halves what the network must learn
and is checked by a unit test on real puzzles.

\begin{itemize}
  \item \textbf{Boards} (2{,}304 bits). Piece-square planes ($12 \times 64$) for three positions:
        before the key move, after it, and after the solver's second move.
  \item \textbf{Moves} (864 bits). Six move blocks: the opponent move that created the position, then
        $m_1, r_1, m_2, r_2, m_3$. Each block holds the from-square, the to-square, the moving piece,
        the captured piece, and flags for check, promotion, under-promotion, castling and en passant.
  \item \textbf{Tactical relations} (90 bits). For each of the first three solver moves: which enemy
        piece types the moved piece attacks, whether it attacks two or more valuable targets, whether
        it is left en prise or defended, discovered and double checks, pins it creates and the types
        of pinned pieces, a piece lined up behind an attacked king, queen or rook, whether the move is
        quiet, a retreat, or captures a defender, and whether the position after it is checkmate.
  \item \textbf{Replies} (6 bits) and \textbf{position facts} (43 bits): forced replies and recaptures;
        hanging pieces on both sides; whether the key move is the only check or the only capture;
        castling rights; the material signature of an endgame; a boxed-in king; the number of solver
        moves in the line; and the mate verdict.
  \item \textbf{Continuous features} (27). Material by piece type and the balance before and after the
        line, the number of legal moves, captures and checks available, pressure on the squares around
        the enemy king and its escape squares, and the number of replies the opponent had.
\end{itemize}

No engine evaluation enters the encoding, so all 5.88M puzzles can be encoded in under an hour on a
laptop, and a position can be encoded during game analysis in about three milliseconds. The one
exception is an optional mate flag, which lets game analysis pass Stockfish's mate verdict for a line
that was cut off before the mate appears.

\subsection{Architecture}

A shared trunk of two fully connected ReLU layers ($3{,}329 \to 512 \to 256$) feeds three linear
heads: a 24-way softmax for the category, 66 independent sigmoids for the themes, and two outputs
$(\mu, s)$ for the difficulty, where $\sigma^2 = e^{s}$. The network has 1.86M parameters. Sharing a
trunk across the three tasks is classical multi-task learning~\cite{caruana1997multitask}: the
representation that says \emph{which} tactic a line contains is largely the representation that says
how hard it is to find.

\subsection{Loss: a noise-aware heteroscedastic objective}

For a batch, with $y_c$ the category, $\mathbf{y}_t$ the theme indicators and $y_b$ the observed
rating (standardised), the loss is

\begin{equation}\label{eq:puzzlenet-loss}
  \mathcal{L} = \underbrace{-\log p_{y_c}}_{\text{category}}
  \;+\; \lambda_t \sum_{j} \mathrm{BCE}\bigl(\sigma(z_{t,j}), y_{t,j}\bigr)
  \;+\; \lambda_r \, \underbrace{\mathrm{sg}\!\left(v^{\beta}\right)}_{\text{$\beta$-NLL weight}}
  \cdot \tfrac{1}{2}\left[\log v + \frac{(y_b - \mu)^2}{v}\right],
\end{equation}

\begin{equation}\label{eq:puzzlenet-var}
  v \;=\; \underbrace{e^{s}}_{\text{content uncertainty}} \;+\; \underbrace{\rho^2}_{\text{known rating noise}} .
\end{equation}

\noindent Two choices in this objective matter for how the model is used later.

\paragraph{Known label noise.} A Lichess rating is not the truth about a puzzle; it is a Glicko-2
estimate that comes with its own deviation $\rho$ (a median of 79 rating points, but far larger for
rarely played puzzles). Equation~\eqref{eq:puzzlenet-var} adds that known measurement variance to the
model's own, so that a noisy label pulls the fit less than a precise one, and so that $e^{s}$ is left
to represent only what the content itself leaves uncertain: how hard this puzzle is, given what is on
the board. That is exactly the quantity the recommender needs for a mined puzzle, which has no
Lichess rating at all, and it is why the variance head is not decoration. Section~\ref{sec:integration-irt}
shows where it enters the Bayesian update.

\paragraph{$\beta$-NLL.} Training a Gaussian likelihood with a learned variance has a known failure
mode: the network can reduce the loss on hard examples by inflating their variance, which in turn
scales down their gradient for the mean, so the hardest examples end up teaching the mean the least.
Seitzer et al.~\cite{seitzer2022pitfalls} correct this by weighting each example's loss by a
stop-gradient copy of $v^\beta$. With $\beta = 1$ the mean receives exactly the plain squared-error
gradient; with $\beta = 0$ the plain likelihood is recovered. This work uses the recommended
$\beta = 0.5$, and both extremes are run as ablations.

\subsection{Optimisation and calibration}

The network, its gradients and its optimiser are written from first principles in
NumPy~\cite{harris2020array}; no deep learning framework is used. Two reasons: the deployed
application keeps no heavy dependency and answers in about a millisecond on a CPU, and every step of
the learning algorithm stays visible and testable. Each hand-derived gradient is checked against
central finite differences in the test suite, for every head and for the linear, single-task and
multi-task configurations.

Training uses AdamW~\cite{loshchilov2019decoupled} (decoupled weight decay on weight matrices only),
gradient-norm clipping, and a learning rate that warms up linearly and then decays on a cosine
schedule~\cite{loshchilov2017sgdr}, with batches of 2{,}048 and He initialisation for the ReLU
layers~\cite{he2015delving}. Model selection keeps the checkpoint with the best validation score.

Raw network outputs are then calibrated on the validation split by fitting two scalars: a softmax
temperature $T$ for the category head~\cite{guo2017calibration}, and a variance scale $c$ such that
$v = c\,e^{s} + \rho^2$ for the difficulty head, in the spirit of post-hoc recalibration for
regression~\cite{kuleshov2018accurate}. Both are one-dimensional optimisations of the validation
likelihood, and both are reported before and after in Chapter~\ref{ch:results}.

\subsection{Data, splits and a fair comparison}

Puzzles are split by a hash of their identifier: 2\% validation, 3\% test, the rest for training.
One further set is carved out first. The rule-based tagger was scored in
Section~\ref{sec:labeller-results} on a held-out sample of 30{,}000 uniformly drawn puzzles plus up to
2{,}000 per category; those 59{,}668 puzzles are \emph{excluded from training and model selection
entirely}, so that both labellers can be scored on exactly the same puzzles, neither having seen them.
The tagger's development sample is not excluded, because the stratified draw takes up to 2{,}000 of
each category and the rarest category has only about 2{,}100 puzzles in total; holding out both
samples would have left the network almost no examples of them.

\subsection{Where the network enters the system}\label{sec:integration-irt}

\begin{enumerate}
  \item \textbf{Labelling critical positions.} Game analysis labels the engine's principal variation
        with PuzzleNet instead of the rule-based tagger. Because game analysis sees an engine line
        rather than a puzzle solution, this shift of input distribution is measured separately
        (Section~\ref{sec:nn-deployment}), and the number of plies the network is shown is chosen from
        that measurement.
  \item \textbf{Rating mined puzzles.} A mined puzzle takes $\mu$ as its rating and $\sigma$ as its
        rating deviation, replacing the closed-form heuristic and the blanket deviation of 300 points.
        The deviation is floored at 75 points, the deviation of a well-played Lichess puzzle: a puzzle
        nobody has ever attempted cannot be known more precisely than that, and the model was
        calibrated on Lichess puzzles rather than on puzzles mined from a player's own games.
  \item \textbf{Feeding the Bayesian recommender.} That per-puzzle deviation is the term $v_b$ in the
        IRT update of Section~\ref{sec:rl-setup}: $\kappa = (1 + 3 v_b / \pi^2)^{-1/2}$ scales how far
        one attempt may move a player's ability and category offsets. An uncertain difficulty estimate
        now moves the player's model less, instead of every mined puzzle being treated as equally
        uncertain.
\end{enumerate}

Every one of these paths falls back to the previous behaviour when no model file is installed or when
\code{PUZZLENET=off} is set, and the rule-based tagger remains available through \code{LABELLER=rules}.
The application therefore still runs, unchanged, without the network.

\subsection{What it is compared against}

The category head is compared with the rule-based tagger on the held-out harness, with paired tests on
the same puzzles (McNemar's exact test~\cite{mcnemar1947note} on per-puzzle correctness and bootstrap
intervals for the differences in $\kappa$ and macro recall). The difficulty head is compared with the
miner's formula, with a constant, with a least-squares fit on line length and the mate verdict, and
with gradient-boosted trees~\cite{ke2017lightgbm} trained on the hand-built features alone. Ablations
trained on a common one-million-puzzle subset isolate the contributions of depth, of the hand-built
features, of the raw board, of multi-task learning and of the $\beta$-NLL exponent, and a
learning-curve sweep from 30{,}000 to 5.5 million puzzles shows how much the data volume matters.
Finally, the labels are carried downstream: into the label-quality simulation, into the 300-player
cohort, and into the prequential replay of the real solve logs.
